\section{Ngomebe --- ``The Stone Men''}
\label{chap:ngomebe}

\subsection*{Quick Reference}
\textbf{Tagline:} \textit{The stone remembers. The ancestors walk with the rollers. The city that stops dies.}

\begin{quote}
  \textbf{Mason-Kin Elder (Elite):}  
  ``A city that stops moving is a city that has chosen to die. The stone remembers every mile, every roller, every oath sworn in the shadow of the keystone. We do not build cities. We raise them from the earth and teach them to walk.''
\end{quote}

\begin{quote}
  \textbf{Stone-Cutter (Commoner):}  
  ``My grandfather carved the keystone for this city. My father greased the rollers. I inherited a sledge and a song. The song is the city's heartbeat. If I stop singing, the city stops walking — and so do I.''
\end{quote}

\begin{quote}
  \textbf{Ancestral Healer (Izinyanga):}  
  ``The ancestors do not speak in words. They speak in cracks. A fissure in the keystone is not a flaw — it is a message. Learn to read the stone's wounds, or the city will read your bones.''
\end{quote}

\vspace{2mm}
\begin{tcolorbox}[colback=black!3,colframe=blue!60!black,title={Theme \& Atmosphere}]
The Ngomebe are a people of deep, iron-dark skin and older ties to the southern ranges. Their cities walk — colossal stone structures on rollers and sledges, dragged by teams of oxen, magic, and Aeler-designed capstans. This partnership is one of the great cross-cultural achievements of the age: Aeler engineers maintain the bearings; Ngomebe masons carve the keystones. A walking city is not a curiosity; it is a living argument against stillness, a monument to survival in a land where staying in one place invites famine, war, or the slow forgetting of ancestral stone.

Spirits are not distant — they live in the rock, the rollers, the cracks. The ancestors speak through keystone fissures. The \textbf{tokoloshe} fouls the bearings at night. The \textbf{impundulu} (lightning bird) shatters entire sledges with a single strike. The Ngomebe do not fight these spirits. They bargain with them: a grindstone of fresh meat, a pot of home-brewed sorghum beer, a promise whispered into a roller-pit.

\vspace{2mm}
\textbf{What you notice first:} The groan of stone on stone. The smell of warm granite, ox sweat, and sorghum beer offerings. The way the ground trembles slightly, even a mile away, as the city inches across the plateau.

\textbf{The one rule that kills newcomers:} Never carve a name into living rock. The stone remembers. It will whisper that name to every spirit of the earth, and the debt will follow your bloodline for seven generations.
\end{tcolorbox}

\subsection*{Regional Mood: The Weight of Movement}

Power here is measured in keystones, migration routes, and the loyalty of the Mason-Kin. A walking city is a polity, a temple, and a weapon all at once. The Mason-Kin Houses control the rights to carve, to move, and to settle. \textbf{Root-binders} — heretics who plant acorns in the rollers — believe the cities should stop and let the forest reclaim the stone. The \textbf{Aeler engineers} are neither masters nor servants; they are partners in a pact older than memory. But beneath them all, the \textbf{Ancestral Healers (Izinyanga)} speak for the dead, and the dead speak through the keystone's cracks. When a city stops, it is because the ancestors have been offended. When a city moves true, it is because the \textbf{Rain Queen} has blessed the path.

\textbf{Starting Location:} The edge of a moving city's path, where the stone is still warm and the dust has not yet settled. A Mason-Kin elder approaches, wiping granite dust from his face: ``A keystone has cracked. The city will stop before dawn. The ancestors are restless. Find the saboteur — a Root-binder, a tokoloshe, or a traitor among us — or push the rollers yourselves. The plateau does not forgive stillness.''

\subsection*{Prominent NPCs of Ngomebe}
\label{sec:ngomebe-npcs}

\begin{longtable}{|p{0.2\textwidth}|p{0.25\textwidth}|p{0.3\textwidth}|p{0.15\textwidth}|}
\hline
\textbf{Name} & \textbf{Role/Title} & \textbf{Motivation} & \textbf{Rumoured Quirk} \\
\hline
Elder Morvath & Mason-Kin elder of the eastern caravan & Keep the city moving south, away from the Red Desert's breath & He has not slept in forty days; the city's rhythm keeps him awake. \\
\hline
Kaelen & Stone-cutter whose city left him behind & Catch up to his family's city before it crosses the Salt Scar & He chisels a new keystone every night. None have been accepted. \\
\hline
Vent-Prior Thrain & Aeler engineer, grease-stained and level-carrying & Maintain the bearings that keep the city alive & She speaks to the rollers in a language of taps. They answer in groans. \\
\hline
The Root-binder & A heretic planting acorns in the roller-pits & Let the forest reclaim the stone; let the cities sleep & He has a single green leaf growing from his left ear. \\
\hline
Ox-Master Senn & Handler of the eight-ox teams that drag the city & Keep the beasts healthy, the yokes oiled, and the pace steady & He names each ox after a dead mason. The oxen remember. \\
\hline
The Still Voice & A spirit of abandoned keystones & Witness the cities that chose to stop & It appears only when a city's migration clock reaches zero. It asks one question: ``Why did you stay?'' \\
\hline
Nomvula (Rain Queen) & Mythical rainmaker, keeper of the drought & Ensure the plateau never dries while the cities walk & She has not been seen in a century. Her staff is a petrified lightning bolt. \\
\hline
Tokoloshe Catcher & Old woman who traps water spirits & Keep the bearings free of malevolent tricksters & She catches tokoloshe in clay pots; they sing when the pot is shaken. \\
\hline
\end{longtable}

\subsection*{Names of Ngomebe}
\label{sec:ngomebe-names}

Inspired by the deep southern plateau: short, earth-toned syllables with clicks and tonal shifts. Surnames often denote a stone type, a city of origin, or a Mason-Kin House.

\begin{longtable}{|p{0.25\textwidth}|p{0.25\textwidth}|p{0.2\textwidth}|p{0.2\textwidth}|}
\hline
\textbf{First Names (Male)} & \textbf{First Names (Female)} & \textbf{Surnames} & \textbf{Epithets} \\
\hline
Morvath & Thandi & Keystone-carver & \textit{the Unstill} \\
\hline
Kaelen & Nia & Roller-born & \textit{the Grey-Faced} \\
\hline
Senn & Zola & Mason-Kin & \textit{the Slow Walker} \\
\hline
Duma & Amara & Stone-song & \textit{the Ninth Roller} \\
\hline
Temba & Lindi & Sledge-heir & \textit{the Root-Breaker} \\
\hline
Jengo & Penda & Ox-kin & \textit{the Warm Stone} \\
\hline
Nomvula & Lethabo & Rain-touched & \textit{the Drought's Daughter} \\
\hline
\end{longtable}

\paragraph*{(Quarry/Hall/Edge)} Abandoned quarry where stone-song still echoes; Mason-Kin hall with walls carved in lineage; the edge of a city's path — the stone is warm to the touch.

\section*{Spades --- Places}
\label{sec:ngomebe-places}
\begin{enumerate}
\setcounter{enumi}{1}
\item \textbf{City on the Move} --- Rollers groan, oxen low, the ground trembles.
  \hfill \textit{The city's shadow moves faster than the city itself. Stay in the shadow to keep pace.}
\item \textbf{Abandoned Quarry} --- Stone-song echoes from the walls.
  \hfill \textit{The song is the memory of the last stone cut here. It sings of a mason who fell.}
\item \textbf{Mason-Kin Hall} --- Walls carved with lineage, floor inlaid with sledge-ruts.
  \hfill \textit{The carvings weep water when a city is in danger. The water tastes of granite.}
\item \textbf{Edge of the Path} --- The stone is warm. The warmth fades when the city leaves.
  \hfill \textit{If you place your palm on the stone before it cools, you may ask the city one question.}
\item \textbf{Roller-Pit} --- A trench where the great rollers turn, greased with rendered ox-fat.
  \hfill \textit{The grease is mixed with ancestor ash. The rollers remember every mason who ever touched them.}
\item \textbf{Keystone Chamber} --- The heart of the city, a single stone that bears all weight.
  \hfill \textit{The keystone hums. The pitch changes when a mason swears a false oath.}
\item \textbf{Sledge-Track Scar} --- A groove worn into the bedrock by centuries of passage.
  \hfill \textit{The groove is deep enough to hide a corpse. Some say it already does.}
\item \textbf{Oxen Corral (Mobile)} --- A pen on sledges, dragged behind the city.
  \hfill \textit{The oxen never sleep. They walk in their sleep. Their dreams move the city.}
\item \textbf{Root-binder's Grove} --- A hidden stand of trees planted in an abandoned roller-pit.
  \hfill \textit{The trees have stone roots. The roots grip the rollers and will not let go.}
\item[J] \textbf{Abandoned City (Dead Keystone)} --- A stone shell on a hill, silent and cold.
  \hfill \textit{The keystone is cracked. A single green shoot grows from the crack. The Root-binder's work.}
\item[Q] \textbf{Bearing Foundry (Aeler)} --- A mobile forge where new rollers are cast.
  \hfill \textit{The forge is pulled by a separate team of oxen. It is always last in the procession.}
\item[K] \textbf{The Thing-Holm of Stone} --- A flat rock where Mason-Kin elders meet.
  \hfill \textit{The rock is a single piece of fallen sky-metal. Oaths sworn here cannot be broken.}
\item[A] \textbf{The First City (Myth)} --- A walking city that never stopped. It is said to circle the plateau, just beyond the horizon.
  \hfill \textit{Those who see it cannot speak of it for a year and a day. The stone seals their lips.}
\end{enumerate}

\paragraph*{(Cutter/Elder/Engineer)} Stone-cutter whose city left them behind; Mason-Kin elder, grey with granite dust; Aeler engineer: grease-stained, carrying a level.

\section*{Hearts --- People \& Factions}
\label{sec:ngomebe-people}
\begin{enumerate}
\setcounter{enumi}{1}
\item \textbf{Stone-Cutter (Left Behind)} --- Her city moved south. She stayed to carve a monument.
  \hfill \textit{She carves her own face into the rock. ``So the city will remember me,'' she says.}
\item \textbf{Mason-Kin Elder} --- Grey with granite dust, eyes the colour of slate.
  \hfill \textit{He carries a hammer that has never struck a false blow. The hammer knows.}
\item \textbf{Root-binder Heretic} --- Planting acorns in the roller-pits, singing to the earth.
  \hfill \textit{His hands are bark. His breath smells of green wood. The oxen fear him.}
\item \textbf{Aeler Engineer} --- Tuning fork in one hand, level in the other.
  \hfill \textit{She can hear a crack in a bearing from a mile away. She never walks under a keystone.}
\item \textbf{Ox-Master} --- Whip of braided leather, voice like a landslide.
  \hfill \textit{He names his oxen after dead masons. The oxen stop when he speaks their names.}
\item \textbf{Ancestral Healer (Izinyanga)} --- Reads the cracks in the keystone, speaks for the dead.
  \hfill \textit{She wears a necklace of ground quartz. The ancestors whisper through the stones.}
\item \textbf{Tokoloshe Catcher} --- Old woman, blind in one eye, carries a clay pot.
  \hfill \textit{She traps tokoloshe in her pots. They sing when shaken. She uses their songs to find sabotage.}
\item \textbf{Impundulu Tamer} --- A legendary bird-whisperer, with scars from lightning.
  \hfill \textit{He can call the lightning bird with a whistle. It demands payment: the blood of a mason.}
\item \textbf{Impi Captain} --- Warrior of the city's flank guard, carrying a shield of braided leather.
  \hfill \textit{She lost two fingers to a Root-binder ambush. She still holds her spear with the remaining three.}
\item[J] \textbf{The First Mason's Ghost} --- A translucent figure in a mason's apron.
  \hfill \textit{He appears when a keystone is first laid. He offers advice. The advice is always wrong.}
\item[Q] \textbf{The Ninth Ox} --- A mythic beast said to pull the First City.
  \hfill \textit{Its hooves never touch the ground. It walks on the shadow of clouds.}
\item[K] \textbf{Bearing-Caster (Aeler)} --- A dwarf who pours molten metal into roller-molds.
  \hfill \textit{He signs each bearing with his own thumbprint. The thumbprint is a prayer.}
\item[A] \textbf{The Still Voice} --- A spirit made of silence and cold stone.
  \hfill \textit{It asks: ``Why did you stop?'' Answer truthfully, or it will never let you leave.}
\end{enumerate}

\paragraph*{(Crack/Sabotage/Tremor)} A keystone cracks — the city stops; Root-binder sabotage — a wheel jams; an earth tremor (all actions Desperate for a scene).

\section*{Clubs --- Complications/Threats}
\label{sec:ngomebe-complications}
\begin{enumerate}
\setcounter{enumi}{1}
\item \textbf{Cracked Keystone} --- The heart of the city fractures. The city stops.
  \hfill \textit{The crack is singing. The song is the name of the mason who carved it. He is long dead.}
\item \textbf{Root-binder Sabotage} --- Acorns sprout in the roller-pits. A wheel jams.
  \hfill \textit{The roots are intelligent. They twist toward the sound of oaths.}
\item \textbf{Earth Tremor} --- The ground shakes. All actions are Desperate for a scene.
  \hfill \textit{The tremor was not natural. The Still Voice is waking.}
\item \textbf{Mason-Kin Feud} --- Two Houses argue over the next migration route.
  \hfill \textit{The argument is carved into the keystone. The stone is splitting.}
\item \textbf{Oxen Plague} --- A sickness spreads through the teams. Speed halves.
  \hfill \textit{The sickness came from a Root-binder's grove. The leaves are in the fodder.}
\item \textbf{Bearing Failure (Aeler)} --- A roller seizes. The engineer blames the stone-cutter.
  \hfill \textit{The bearing was sabotaged with sand from the Red Desert. The sand whispers.}
\item \textbf{Sledge-Track Collapse} --- The ground gives way. A section of the city sinks.
  \hfill \textit{The collapse reveals an older city buried beneath. The older city is still walking.}
\item \textbf{Migration Blockade} --- Another city has stopped on the planned route.
  \hfill \textit{The stopped city's keystone is dead. Its people have become Root-binders.}
\item \textbf{Ancestral Debt Called} --- A Mason-Kin House demands a sledge-right from a century ago.
  \hfill \textit{The debt is carved on a tablet of slate. The slate is warm. The ancestors are watching.}
\item[J] \textbf{Tokoloshe in the Bearings} --- A water spirit fouls the roller-pits. The city groans.
  \hfill \textit{The tokoloshe is small, fast, and invisible to the unaware. Only a catcher can see it.}
\item[Q] \textbf{Impundulu Strikes} --- A lightning bird shatters a keystone with a thunderous blow.
  \hfill \textit{The bird is angry. Someone broke a vow to the Rain Queen. The city bleeds stone.}
\item[K] \textbf{Root-binder Revolt} --- A grove of walking trees attacks the roller-pits.
  \hfill \textit{The trees have stone bark. The stones are carved with Mason-Kin faces.}
\item[A] \textbf{The First City Approaches} --- A city older than memory appears on the horizon.
  \hfill \textit{It is made of obsidian. Its rollers are giant femurs. It wants to merge.}
\end{enumerate}

\paragraph*{(Charter/Wax/Song)} Keystone charter (right to dwell in a city for a season); sledge-wax (greases negotiation with the Mason-Kin); stone-song fragment (one reroll on a survival test).

\section*{Diamonds --- Rewards/Leverage}
\label{sec:ngomebe-rewards}
\begin{enumerate}
\setcounter{enumi}{1}
\item \textbf{Keystone Charter} --- Right to dwell in a city for a season, tax-free.
  \hfill \textit{The charter is a chip of the keystone itself. Carry it to be welcomed.}
\item \textbf{Sledge-Wax} --- A lump of grey grease. Use it to gain +1 Position in Mason-Kin negotiations.
  \hfill \textit{The wax is rendered from the fat of an ox that walked a thousand miles.}
\item \textbf{Stone-Song Fragment} --- A recorded hum. Reroll a failed Survival test once.
  \hfill \textit{The song is the city's heartbeat. Hum it, and the stone will guide you.}
\item \textbf{Roller-Bearing (Aeler-crafted)} --- A spare bearing. Gift it to an engineer for a favour.
  \hfill \textit{The bearing is stamped with a rune that means ``no debt.'' The rune is a lie.}
\item \textbf{Mason-Kin Hammer} --- A tool that never misses a keystone's true face.
  \hfill \textit{The hammer's head is a piece of the First City's keystone. It whispers warnings.}
\item \textbf{Ox-Master's Whip} --- Crack it to make oxen move faster for one leg.
  \hfill \textit{The whip's braid includes a hair from every ox that ever pulled a city.}
\item \textbf{Root-binder's Acorn} --- Plant it to stop a city's roller for one day.
  \hfill \textit{The acorn grows into a sapling within an hour. The sapling has stone roots.}
\item \textbf{Abandoned City Map} --- A chart of dead cities and their keystone graves.
  \hfill \textit{The map is drawn on stone. The stone is a page. The pages are bound in granite.}
\item \textbf{Elder's Verdict} --- A ruling that settles a Mason-Kin dispute in your favour.
  \hfill \textit{The verdict is carved on a slate tablet. Carry it as a shield.}
\item[J] \textbf{Bearing-Foundry Pass} --- Access to the Aeler mobile forge for one repair.
  \hfill \textit{The pass is a single steel coin. The coin is still hot.}
\item[Q] \textbf{Still Voice's Silence} --- One question answered truthfully by the spirit of abandoned stone.
  \hfill \textit{The answer comes as a cold breath. The breath smells of geode crystals.}
\item[K] \textbf{First City's Echo} --- A piece of obsidian that hums with ancient rhythm.
  \hfill \textit{Once per campaign, touch it to make a city move one mile in a single night.}
\item[A] \textbf{The Living Keystone} --- A fragment from a city that still walks.
  \hfill \textit{Plant it in dead stone, and a new city will grow. The growth takes a century.}
\end{enumerate}

\section*{Quick use notes}
\label{sec:ngomebe-quick-use}
\begin{itemize}
\item Draw until you have all four suits: Spade = place, Heart = actor, Club = pressure, Diamond = leverage. Highest rank sets the main Clock (2–5 → 4, 6–10 → 6, J/Q/K → 8, A → 10).
\item Diamonds are codified outcomes (charters, wax, songs) that change position rather than call for a roll.
\item If any A appears, echo \textbf{stone-omens}—rollers that groan in harmony, keystones that weep water, a city that leaves no track.
\end{itemize}

\subsection*{Additional Features (Cultural Mechanics)}
\label{sec:ngomebe-mechanics}

\begin{itemize}
\item \textbf{The Keystone Bond (Relationship Mechanic):} Each walking city has a Keystone — a single stone that bears the weight of the roof. The Keystone is also the physical anchor of the city's ancestral spirit. A character who bonds with the Keystone (by spending a downtime scene in meditation at its base) gains:
  \begin{itemize}
    \item \textbf{Keystone Sense:} Once per scene, the character can ask the Keystone one question about the city's current state (``Is the eastern roller cracking?'' ``Is there a saboteur in the foundry?''). The Keystone answers truthfully, but only about the present.
    \item \textbf{Keystone Debt:} The character gains a Debt Clock (size 4) to the ancestor. When it fills, the ancestor demands a service — usually a pilgrimage to a forgotten quarry or the retrieval of a lost roller.
  \end{itemize}
\item \textbf{The Roller-Pits (Hazard Mechanic):} The city moves on rollers — massive wooden cylinders that must be greased and rotated. The roller-pits are dangerous, loud, and prone to sabotage.
  \begin{itemize}
    \item Entering a roller-pit requires a \textbf{Body + Athletics test (DV 4)} to avoid being crushed. On a failure, the character takes 1 Harm.
    \item Sabotaging a roller-pit (to stop the city or redirect it) requires a \textbf{Wits + Craft test (DV 5)}. On a success, the city moves in the desired direction; on a failure, the roller jams and the city stops — and the ancestors are angry.
    \item Repairing a jammed roller takes 3 actions. Each action requires a \textbf{Body + Craft test (DV 4)}. On a failure, the roller shifts and the character takes 1 Fatigue.
  \end{itemize}
\item \textbf{The Migration Clock (Campaign Clock, size 10):} Every walking city has a Migration Clock. It advances each day the city moves, and retreats when the city stops.
  \begin{itemize}
    \item At 2: The city enters fertile land. The party gains +1 die to foraging and hunting.
    \item At 4: The city passes a rival settlement. The party may trade or raid.
    \item At 6: The city enters a narrow pass. All navigation rolls are Desperate.
    \item At 8: The city approaches a sacred site. The ancestors grow restless — all rolls involving the Keystone gain +1 die, but failure ticks the Keystone Debt clock by 2.
    \item At 10: The city reaches its destination — the end of the migration. The ancestors are satisfied. The party may claim a \textbf{Keystone Blessing} (a permanent +1 die to all rolls involving stonework or negotiation with Ngomebe).
  \end{itemize}
  The party can influence the Migration Clock by:
  \begin{itemize}
    \item \textbf{Scouting ahead} (Wits + Survival) – success advances the clock by 1, failure ticks it back by 1.
    \item \textbf{Repairing rollers} (Craft test) – success advances the clock by 2, failure advances it by 0 but ticks the Keystone Debt.
    \item \textbf{Paying tribute} (a Honey Jar, a story, or a Boon) – advances the clock by 1 without risk.
  \end{itemize}
\item \textbf{The Root-binder Heresy (Faction Mechanic):} Root-binders are heretics who plant acorns in the roller-pits, hoping to stop the cities and let the forest reclaim the stone. They have a \textbf{Sabotage Clock (size 6)}. Each time the party fails a Craft or Survival roll in a roller-pit, the Sabotage Clock ticks.
  \begin{itemize}
    \item At 2: A roller cracks. The city loses 1 Speed (travel takes longer).
    \item At 4: A roller-pit floods. The city must stop for repairs (lose 1 day).
    \item At 6: A major collapse. The city cannot move until the party clears the pit (requires a Strength + Craft test, DV 5, and costs 1 day).
  \end{itemize}
  Players can work with the Root-binders (gain information, sabotage rival cities) or against them (hunt them, protect the rollers). The Root-binders offer \textbf{Green Blessings} – a permanent +1 die to survival in the forest, but at the cost of a Debt Clock to ``the forest itself'' (the Hollow's attention).
\item \textbf{The Mason-Kin Houses:} The walking cities are not nations but extended families. Each House controls a keystone, a migration route, and the rights to carve. Houses feud over sledge-tracks, water rights, and the honour of the First City.
\item \textbf{The Aeler Pact:} The Aeler engineers of the southern ranges maintain the bearings and capstans that make the walking cities possible. The Ngomebe provide keystone hearts and stone-song. It is a partnership of equals, sealed in slate and iron.
\item \textbf{Ancestral Cracks:} Once per session, a character may press their ear to a keystone and roll Spirit + Lore (DV 3). On success, the ancestors whisper a warning or a clue (GM reveals a hidden complication or a shortcut). On a miss, the stone is silent — but the ancestors are listening.
\item \textbf{Tokoloshe Pots:} A character who carries a clay pot with a tokoloshe inside gains +1 die to detect sabotage or traps in the city. The tokoloshe sings when danger is near. However, if the pot breaks, the tokoloshe escapes and will haunt the carrier for a full season (all actions Desperate at night).
\item \textbf{Local Patrons (Syncretic Names):}
  \begin{itemize}
    \item \textbf{The First Stone} — The Sealed Gate. *The keystone that never cracks; the oath that never breaks.*
    \item \textbf{The Rain Queen} — Raéyn. *The bringer of storms and the drought; the mistress of the path.*
    \item \textbf{The Lightning Bird} — The Storm King. *The shatterer of false keystones; the judge of masons.*
    \item \textbf{The Ancestor's Voice} — The Witness. *The whisper in the crack; the memory in the stone.*
  \end{itemize}
\end{itemize}

\subsection*{Lore Echoes (Cross-Expansion Hooks)}
\begin{itemize}
  \item \textbf{The Red Desert's Breath (Ashaan):} The Red Desert is spreading. Some Ngomebe cities are moving south to escape the dust. The Keystone Charter of a southern city may be the only thing standing between a tribe and annihilation.
  \item \textbf{The Sledge-Wax Trade (Way of Silk):} Sledge-wax is a valuable commodity on the caravan routes. Tulkani merchants trade it for story-knots; Fhara water-clerks use it to seal well-contracts. A missing shipment of sledge-wax could stop a city.
  \item \textbf{The Iron Mines of Vajrapur (Dhahara):} The Aeler engineers of Ngomebe source their bearing-metal from the Dhaharan iron mines. A dispute between a Mason-Kin House and a Dhaharan raja could choke the supply of rollers.
  \item \textbf{The Hollow and the Still City (Eastern Realms):} The Hollow is attracted to still places. A city that stops moving becomes a beacon. The Aelaerem whisper that the First City never stopped because it is running from something.
  \item \textbf{Tokoloshe and the Bitter Root:} Daughters of the Cord sometimes capture tokoloshe and weave them into Hollowed children. The tokoloshe inside the child sings when a debt is due. The Ngomebe avoid Daughter-run hospices for this reason.
\end{itemize}

\subsection*{Ngomebe Superstitions (burn for +1 die once)}
\begin{itemize}
  \item \textbf{Bread to the Roller:} Crumble a crust into the grease-pit. Ignore the first \emph{Earth Tremor} penalty. `[SUPERSTITION]`
  \item \textbf{Knot the Sledge-Rope:} Tie a single knot in an ox's halter. Cancel \emph{Oxen Plague} or take –1 die on your next handling roll (your choice). `[SUPERSTITION]`
  \item \textbf{Name to the Keystone:} Whisper a dead mason's name into the keystone's crack. Gain +1 Effect when negotiating with the Mason-Kin this scene, but mark \emph{Obligation} to the First City. `[SUPERSTITION]`
\end{itemize}

\subsection*{Thematic SB Spend Table}
\label{sec:ngomebe-sb}

\paragraph{Minor Complications (1 SB)}
\begin{itemize}
\item \textbf{Exposure:} Your actions draw unwanted attention from \textbf{Mason-Kin guards or Aeler engineers}.
\item \textbf{Noise:} Sounds of your actions alert nearby \textbf{ox-drivers or stone-cutters}.
\item \textbf{Trace:} Evidence of your passage marks your route for \textbf{Root-binders or keystone guardians}.
\item \textbf{Delay:} A brief but meaningful setback costs you \textbf{time or favorable migration pace}.
\item \textbf{Supply Strain:} Mark +1 segment on a relevant \textbf{resource clock}.
\end{itemize}

\paragraph{Moderate Setbacks (2 SB)}
\begin{itemize}
\item \textbf{Alarm Raised:} \textbf{Local Mason-Kin elder or Aeler vent-prior} becomes aware and begins responding.
\item \textbf{Position Lost:} You lose advantageous ground/cover/stealth due to \textbf{earth tremor or roller jam}.
\item \textbf{Foe Appears:} A \textbf{Root-binder cell or rival House guard} arrives on scene.
\item \textbf{Gear Trouble:} A piece of equipment becomes \textbf{Compromised/Neglected}.
\item \textbf{Lock/Barrier:} A simple obstacle now requires a test to overcome.
\end{itemize}

\paragraph{Serious Trouble (3 SB)}
\begin{itemize}
\item \textbf{Reinforcements:} Additional \textbf{Mason-Kin, Root-binders, or Aeler engineers} arrive.
\item \textbf{Key Gear Breaks:} A crucial tool/weapon becomes temporarily unusable.
\item \textbf{Major Twist:} The situation fundamentally changes—\textbf{keystone cracks/city stops/First City sighted}.
\item \textbf{Rail Tick:} Advance a relevant campaign/front clock by 1 segment.
\item \textbf{Condition Applied:} Mark \textbf{Fatigue 1/Harm 1/Condition} appropriate to fiction.
\end{itemize}

\paragraph{Major Turns (4+ SB)}
\begin{itemize}
\item \textbf{Trap Springs:} A prepared danger activates with full effect.
\item \textbf{Authority Arrival:} \textbf{Mason-Kin High Council, Still Voice, or First City's herald} intervenes.
\item \textbf{Scene Shift:} The environment changes dramatically—\textbf{tremor opens a chasm/city splits in two/keystone begins to sing}.
\item \textbf{Patron Omen:} Divine/arcane forces take notice—\textbf{omen appears/blessing lost/curse manifests}.
\item \textbf{Narrative Pivot:} The story takes an unexpected turn that reframes objectives.
\end{itemize}

\paragraph{Region-Specific SB Options}
\begin{itemize}
\item \textbf{Ngomebe (Stone Law):} Keystone hums a warning, roller-grooves glow faintly, the ground trembles without quake.
\item \textbf{Ngomebe (Root-binder Heresy):} Acorns sprout from cracks, vines write warnings in the dust, trees move against the wind.
\item \textbf{Ngomebe (Aeler Pact):} Bearings click in counterpoint, tuning forks ring without being struck, the smell of hot metal hangs in still air.
\item \textbf{Ngomebe (Ancestral Spirits):} A cold breath whispers a name, the keystone weeps water, a tokoloshe laugh echoes from a dry roller-pit.
\end{itemize}

\subsection*{Pursuit Ladder (Near / Mid / Far)}
\begin{itemize}
\item \textbf{Advance/Withdraw (Wits + Survival):} On a hit, shift one band. On a strong hit, also impose \emph{False Trail}.
\item \textbf{Cut the Sledge (Presence + Sway):} On a hit, freeze range for one exchange; if a keystone charter is in your possession, +1 die.
\item \textbf{Weather the Tremor (Body + Athletics):} On a hit in \emph{Earth Tremor}, you pick who stays standing; on a miss, the GM does.
\end{itemize}

\subsection*{Boss Seeds (Stone \& Root)}
\begin{itemize}
\item \textbf{The Root-King (Nine Groves)} — \emph{Phases}: Acorn Sabotage $\rightarrow$ Orchestrated Stillness $\rightarrow$ Forest Coup. \emph{Cracks}: burn a grove, expose a false acorn, survive a landslide.
\item \textbf{The Still Voice (Silence of Stone)} — \emph{Phases}: Abandoned Keystones $\rightarrow$ Procession of Dead Cities $\rightarrow$ Stillness Baptism of an Elder. \emph{Cracks}: relic stone-song, survivor testimony, the city moving against the Voice.
\item \textbf{The Tokoloshe King (Master of Tricksters)} — \emph{Phases}: Bearings Foul $\rightarrow$ Roller Pit Flooded $\rightarrow$ Oxen Stampede. \emph{Cracks}: trap him in a clay pot, negotiate with a pot of sorghum beer, feed him a false name.
\end{itemize}

\begin{tcolorbox}[colback=black!3,colframe=black!40!white,title={Patronage \& Power}]
Among the Ngomebe, power flows through keystones, migration rights, and the loyalty of the Mason-Kin Houses. An elder's authority rests on the age of his House's keystone and the number of cities that follow his route. The Aeler engineers hold technical authority; the Root-binders offer a heretical alternative; and the Ancestral Healers speak for the dead, who still vote through the cracks in the stone.

\textbf{For the GM:}  
Power in Ngomebe society is measured in movement and stone. Patronage often takes the form of keystone charters, sledge-wax, or stone-song fragments—each tied to a specific city and House. To emphasize this:
\begin{itemize}
\item Tie rewards to physical tokens (chips of keystone, lumps of wax, recorded songs) that can be challenged, stolen, or voided.
\item Let rival Houses issue conflicting claims, forcing players to choose whose city shelters them.
\item Use the roller-pits, keystone chambers, and Mason-Kin halls as arenas for social contests, where migration routes are debated and saboteurs are judged.
\item Introduce \textbf{Ancestral Cracks} as a way for players to gain clues from the dead — but at the cost of the ancestors' attention.
\end{itemize}
In Ngomebe society, the stone remembers everything — including your debts. The ancestors do not forget. The tokoloshe does not forgive. And the city that stops dies.
\end{tcolorbox}

% Index entries
\index{Ngomebe|see{Stone Men}}
\index{Theme generator}
\index{Stone Men generator}
\index{Places!Ngomebe}
\index{People!Ngomebe}
\index{Complications!Ngomebe}
\index{Rewards!Ngomebe}
\index{Mason-Kin}
\index{Root-binder}
\index{Aeler engineers}
\index{Keystone}
\index{Walking city}
\index{First City}
\index{Still Voice}
\index{Ox-Master}
\index{Stone-song}
\index{Sledge-wax}
\index{Roller-bearing}
\index{Ancestral Healers}
\index{Izinyanga}
\index{Tokoloshe}
\index{Impundulu}
\index{Rain Queen}
\index{Local Patrons!Ngomebe}